\documentclass[11pt,letterpaper]{article}

\usepackage[utf8]{inputenc}
\usepackage[spanish]{babel}
\usepackage[margin=2.5cm]{geometry}
\usepackage{amsmath,amssymb}
\usepackage{enumitem}
\usepackage{titlesec}
\usepackage{array}
\usepackage{hyperref}
\usepackage{xcolor}

\definecolor{unicaucaverde}{RGB}{0,90,60}
\hypersetup{colorlinks=true,linkcolor=unicaucaverde,urlcolor=unicaucaverde}

\titleformat{\section}{\normalfont\Large\bfseries\color{unicaucaverde}}{\thesection}{1em}{}
\titleformat{\subsection}{\normalfont\large\bfseries}{\thesubsection}{1em}{}

\setlength{\parindent}{0pt}
\setlength{\parskip}{6pt}

\title{\textcolor{unicaucaverde}{\Large Guía de Aprendizaje — Semana 1}\\[4pt]
  \large Sistemas de ecuaciones lineales: introducción y el caso de dos incógnitas}
\author{Álgebra Lineal para Ingeniería --- Universidad del Cauca}
\date{2026}

\begin{document}
\maketitle

\begin{center}
\begin{tabular}{|>{\bfseries}p{4cm}|p{10cm}|}
\hline
Semana & 1 de 16 \\ \hline
Unidad & 1 --- Sistemas de ecuaciones lineales (Grossman, cap. 1) \\ \hline
Subtemas & 1.1 (introducción), 1.2 (dos ecuaciones con dos incógnitas) \\ \hline
Horas de clase & 4 \\ \hline
Resultado de aprendizaje & RA1 (parte 1 de 2; ver \texttt{guia1/guia.tex}) \\ \hline
\end{tabular}
\end{center}

\section*{Relevancia para ingeniería}

Un sistema de ecuaciones lineales aparece cada vez que varias cantidades desconocidas están
ligadas por relaciones lineales simultáneas: corrientes y voltajes en un circuito (leyes de
Kirchhoff), fuerzas en equilibrio en una armadura, balances de masa en un proceso, o precios de
equilibrio oferta-demanda. Resolver el sistema es encontrar los valores que satisfacen todas las
relaciones a la vez. El caso de dos ecuaciones con dos incógnitas ya muestra los tres
comportamientos posibles (solución única, ninguna, o infinitas) que se generalizan al caso de
$m$ ecuaciones con $n$ incógnitas (Semana 2).

\section{Objetivo de la semana}

Reconocer la forma general de un sistema de ecuaciones lineales, y resolver e interpretar
geométricamente sistemas de dos ecuaciones con dos incógnitas.

\section{Resultados de aprendizaje de la semana}

\begin{itemize}[leftmargin=1.2cm]
  \item[\textbf{RA1.1.}] Escribir la forma general de un sistema de $m$ ecuaciones lineales con
        $n$ incógnitas, identificar sus coeficientes, y determinar si una tupla dada es solución.
  \item[\textbf{RA1.2.}] Resolver un sistema de dos ecuaciones con dos incógnitas por sustitución
        y por eliminación, clasificarlo (solución única, sin solución, infinitas soluciones), e
        interpretar cada caso como la intersección de dos rectas.
\end{itemize}

\section{Contenidos}

\begin{itemize}[leftmargin=1.2cm]
  \item \textbf{1.1.} Introducción a los sistemas de ecuaciones lineales.
  \item \textbf{1.2.} Dos ecuaciones lineales con dos incógnitas.
\end{itemize}

\section{Plan de la sesión (4 horas)}

\begin{enumerate}[leftmargin=1.2cm]
  \item[\textbf{Hora 1.}] Ecuación lineal y sistema de $m$ ecuaciones con $n$ incógnitas: forma
        general, coeficientes, término independiente, solución, sistema consistente/inconsistente
        (teoría + ejemplos).
  \item[\textbf{Hora 2.}] Taller: verificar si una tupla dada es solución de un sistema;
        traducir un problema de ingeniería sencillo (circuito, mezcla) a un sistema de
        ecuaciones.
  \item[\textbf{Hora 3.}] Dos ecuaciones con dos incógnitas: solución por sustitución y por
        eliminación; interpretación geométrica como dos rectas (teoría + ejemplos de los tres
        casos).
  \item[\textbf{Hora 4.}] Taller: resolver y clasificar sistemas $2\times 2$; exploración
        interactiva en \texttt{visualizacion/} (sección \emph{1.1--1.2 Dos ecuaciones con dos
        incógnitas}) variando los coeficientes para ver los tres casos.
\end{enumerate}

\section{Actividades}

\begin{itemize}[leftmargin=1.2cm]
  \item \textbf{En clase}: talleres de las horas 2 y 4 (ver arriba).
  \item \textbf{Trabajo autónomo}: lista de ejercicios de \texttt{notas.tex} §1--2, de entrega
        antes de la Semana 2.
  \item \textbf{Software}: en \texttt{visualizacion/index.html}, sección \emph{1.1--1.2}, variar
        los coeficientes $a,b,c$ de dos rectas $a_1x+b_1y=c_1$, $a_2x+b_2y=c_2$ y observar cuándo
        el sistema tiene solución única, ninguna, o infinitas.
\end{itemize}

\section{Material de apoyo}

\begin{itemize}[leftmargin=1.2cm]
  \item \texttt{notas.tex} --- desarrollo teórico completo de 1.1 y 1.2, con ejemplos y ejercicios.
  \item \texttt{diapositivas.tex} --- síntesis visual de la sesión.
  \item \texttt{visualizacion/index.html} --- sección \emph{1.1--1.2 Dos ecuaciones con dos
        incógnitas} (interactivo).
\end{itemize}

\section{Evaluación de la semana}

Control formativo corto al inicio de la Semana 2 (10 minutos): 2--3 sistemas $2\times 2$ para
resolver y clasificar, para verificar RA1.1 y RA1.2 antes de avanzar a 1.3--1.4. Contribuye al
Quiz de la Unidad 1 (ver \texttt{guia1/guia.tex} § Evaluación).

\section{Bibliografía de la semana}

\begin{enumerate}[leftmargin=1.2cm]
  \item Grossman, S. I., Flores Godoy, J. J. \emph{Álgebra lineal}. 7a/8a ed., McGraw-Hill,
        cap. 1, §1.1--1.2.
  \item Lay, D. C., Lay, S. R., McDonald, J. J. \emph{Álgebra lineal y sus aplicaciones}. Pearson.
\end{enumerate}

\end{document}
